\section{Metodología de Trabajo}
Para la gestión del ciclo de vida del desarrollo se adoptará un enfoque metodológico híbrido que combina prácticas de \textbf{Scrum} y \textbf{Kanban}, asegurando adaptabilidad ante cambios de requisitos e iteraciones rápidas sin perder la predictibilidad del alcance global.

\subsection{Definición de Roles}
\begin{itemize}
    \item \textbf{Equipo de Alumnos:} Responsables directos de la arquitectura, codificación, aseguramiento de la calidad (QA), descubrimiento de requisitos, mitigación de riesgos y documentación técnica.
    \item \textbf{Cuerpo de Tutores:} Actuarán como facilitadores y \textit{Stakeholders} clave, validando el incremento de valor al cierre de cada ciclo y asistiendo en la priorización del \textit{Product Backlog}.
\end{itemize}

\subsection{Dinámica de Trabajo y Gestión de Tiempos}
\begin{itemize}
    \item \textbf{Sprints Bisemanales:} Ciclos fijos de desarrollo orientados a la entrega de módulos potencialmente desplegables.
    \item \textbf{Sprint Planning y Estimación:} Planificación al inicio de cada sprint utilizando la técnica de \textit{Planning Poker} basada en la escala de Fibonacci para asignar puntos de historia de manera consensuada.
    \item \textbf{Review y Retrospectiva:} Reuniones al cierre de cada iteración para demostrar las funcionalidades ante los tutores, recibir feedback y refinar los procesos internos del equipo.
\end{itemize}